\begin{table}[!htbp]\centering\scriptsize\setlength{\tabcolsep}{4pt}
\caption{Ce que paieraient des ménages types (IR + CSG + CRDS, € par an)}\label{tab:ircsg_cas}
\begin{tabular}{@{}lrrr@{}}
\toprule
& Aujourd'hui & \multicolumn{2}{c}{Variation} \\ \cmidrule(l){3-4}
& & Réforme naïve & Proposition \\
\midrule
Jeune salarié à mi-temps au Smic & 1\,014 & $-$645 & $-$812 \\
Jeune salarié au Smic & 2\,027 & $-$480 & $-$447 \\
Cadre seul, 60 000 € brut & 12\,190 & $-$954 & $-$2\,354 \\
Cadre seul, 100 000 € brut & 24\,973 & $-$976 & $-$4\,015 \\
Cadre dirigeant seul, 200 000 € brut & 64\,992 & $-$695 & $-$10\,917 \\
Parent isolé, deux enfants, au Smic & 2\,027 & $-$1\,290 & $-$2\,027 \\
Couple sans enfant, 50 000 € brut (60/40) & 5\,517 & $-$1\,436 & $-$1\,328 \\
Couple sans enfant à un seul revenu, 80 000 € & 11\,747 & $-$1\,260 & +3\,498 \\
Couple avec deux enfants, 50 000 € brut (60/40) & 4\,765 & $-$1\,949 & $-$2\,408 \\
Couple avec deux enfants à un seul revenu, 50 000 € & 4\,765 & $-$1\,952 & $-$862 \\
Couple avec trois enfants, 60 000 € brut (60/40) & 5\,718 & $-$2\,844 & $-$2\,922 \\
Couple avec trois enfants à un seul revenu, 60 000 € & 5\,718 & $-$2\,832 & $-$1\,003 \\
Retraité seul, pension de 12 000 € & 0 & +524 & 0 \\
Retraité seul, pension de 16 000 € & 0 & +1\,120 & +1\,005 \\
Retraité seul, pension de 20 000 € & 860 & +857 & +1\,352 \\
Retraité seul, pension de 25 000 € & 2\,270 & +192 & +1\,815 \\
Retraité seul, pension de 35 000 € & 5\,138 & $-$105 & +1\,732 \\
Couple de retraités, 30 000 € de pensions (60/40) & 1\,710 & +233 & $-$109 \\
Couple de retraités, 45 000 € de pensions (60/40) & 4\,905 & $-$726 & +1\,330 \\
\bottomrule
\end{tabular}
\par\vspace{3pt}\parbox{\linewidth}{\raggedright\scriptsize\color{encre}
Cas-types OpenFisca-France (législation des revenus 2024) : salariés du privé et retraités de plus de 65 ans, sans autre revenu ni niche ; revenus bruts ; couples mariés ou pacsés, revenus partagés 60/40 sauf « un seul revenu » ; Smic 2024 de 21\,270~€ par an. Réforme naïve : barème de l'IR du foyer (quotient conjugal et familial) et taux proportionnel de 3,91~\% sur l'assiette large. Proposition : sous-partie \ref{sec:ircsg}.4. Retraité seul : sa charge d'aujourd'hui dépend du revenu fiscal de l'avant-dernière année (régime stable supposé).}
\end{table}
